% =========================================================
% Accumulation: Integration — Notes Index
% =========================================================

% ---------------------------------------------------------
% Breadcrumb
% ---------------------------------------------------------
\begin{tcolorbox}[
  colback=gray!6,
  colframe=gray!40,
  arc=2pt,
  left=8pt, right=8pt, top=6pt, bottom=6pt,
  title={\small\textbf{Where You Are in the Journey}},
  fonttitle=\small\bfseries
]
\begin{center}
\small
$\mathbb{N}$, $\mathbb{Z}$, $\mathbb{Q}$, $\mathbb{R}$
$\;\to\;$ Real Analysis
$\;\to\;$ Finiteness (Compactness)
$\;\to\;$ Control (Continuity)
$\;\to\;$ \textbf{Accumulation (Integration)}
$\;\to\;$ Measuring (Measure Theory)
$\;\to\;$ $\cdots$
\end{center}

\medskip
\noindent\textbf{How we got here.}
Continuous functions on compact sets are uniformly continuous.
Uniform continuity is precisely the condition that makes Riemann
sums partition-independent: no matter how we subdivide the
interval, the upper and lower sums are forced to converge to the
same value.
The long structural climb from field axioms through sequences,
topology, compactness, and continuity has been building toward
this payoff.

\medskip
\noindent\textbf{What this chapter builds.}
The Riemann integral via Darboux upper and lower sums, the
Riemann integrability criterion ($\sup L = \inf U$), the
integrability of continuous and monotone functions, linearity
and monotonicity of the integral, the Mean Value Theorem for
integrals, and the Fundamental Theorem of Calculus in both
directions.
The chapter closes by identifying the limits of Riemann
integration --- the functions it cannot handle --- which
motivates the passage to measure theory.

\medskip
\noindent\textbf{Where this leads.}
The Riemann integral works beautifully for continuous and piecewise
continuous functions, but fails for functions with too many
discontinuities or defined on unbounded domains.
Measure theory rebuilds integration from the ground up, replacing
interval-based area with an abstract notion of size on arbitrary
measurable sets.
The Lebesgue integral extends the Riemann integral and admits
the convergence theorems --- monotone convergence, dominated
convergence, Fatou's lemma --- that make rigorous probability
theory and functional analysis possible.
\end{tcolorbox}
\vspace{1em}

% ---------------------------------------------------------
% Structural Role
% ---------------------------------------------------------
\begin{tcolorbox}[
  colback=gray!6,
  colframe=gray!40,
  arc=2pt,
  left=8pt, right=8pt, top=6pt, bottom=6pt,
  title={\small\textbf{Structural Role: Accumulation — Infinite Sums with Finite Precision}},
  fonttitle=\small\bfseries
]
Integration is the study of accumulation.

The central idea is that the area under a curve can be approximated
by finite sums of rectangular slices, and that as the slices become
arbitrarily fine, these sums converge to a definite value.
The entire machinery of this project --- completeness,
sequences, uniform continuity --- is what guarantees that
convergence actually occurs.

The Darboux approach makes this precise.
Upper sums overestimate; lower sums underestimate.
Integrability means the overestimates and underestimates are forced
to agree in the limit.
Uniform continuity on compact domains is what closes the gap.

The Fundamental Theorem of Calculus then reveals the deepest
structural connection in elementary analysis:
\[
  \text{Differentiation} \;\longleftrightarrow\; \text{Integration.}
\]
These are not merely inverse operations.
The FTC says that the accumulated change of a rate equals the net
change of the quantity --- a statement about the relationship
between local behavior (the derivative) and global accumulation
(the integral).

Integration is the culmination of real analysis on $\mathbb{R}$.
Everything that follows --- measure theory, $L^p$ spaces,
Fourier analysis --- is either a generalization or a consequence.
\end{tcolorbox}
\vspace{1em}

% ---------------------------------------------------------
% Structural Roadmap
% ---------------------------------------------------------
\subsection*{Structural Roadmap}

\begin{center}
\textbf{Definitions $\longrightarrow$ Main Theorems
$\longrightarrow$ Consequences and Logical Structure}
\end{center}

The development builds the Riemann integral from partitions
upward, using the completeness of~$\mathbb{R}$ and the
uniform continuity machinery from the previous chapter.
The global progression is:

\begin{enumerate}
  \item Partitions, tagged partitions, and Riemann sums
  \item Darboux upper and lower sums
  \item The Riemann integrability criterion
  \item Integrability of continuous functions (uses uniform
        continuity on compact sets)
  \item Integrability of monotone functions
  \item Linearity, monotonicity, and additivity of the integral
  \item Mean Value Theorem for integrals
  \item Fundamental Theorem of Calculus: Part I (antiderivative)
        and Part II (evaluation)
  \item Improper integrals and the limits of Riemann integration
  \item Transition: what the Riemann integral cannot do, and why
        Lebesgue measure is needed
\end{enumerate}

\begin{remark*}[Primary sources]
Abbott, \textit{Understanding Analysis}, Chapter~7;
Pugh, \textit{Real Mathematical Analysis}, Chapter~4;
Rudin, \textit{Principles of Mathematical Analysis}, Chapter~6.
\end{remark*}

% ---------------------------------------------------------
% Status
% ---------------------------------------------------------
\vspace{1em}
\begin{tcolorbox}[
  colback=gray!6, colframe=gray!40, arc=2pt,
  left=8pt, right=8pt, top=6pt, bottom=6pt,
  title={\small\textbf{Status: Planned}},
  fonttitle=\small\bfseries
]
\begin{center}\Large\bfseries Coming Soon\end{center}
\vspace{6pt}
\noindent Notes, proofs, and exercises will appear here in a future revision.
\end{tcolorbox}

% Future sections (uncomment as content is written):
% \input{volume-iii/analysis/integration/notes/notes-partitions}
% \input{volume-iii/analysis/integration/notes/notes-darboux}
% \input{volume-iii/analysis/integration/notes/notes-integrability}
% \input{volume-iii/analysis/integration/notes/notes-properties}
% \input{volume-iii/analysis/integration/notes/notes-ftc}
% \input{volume-iii/analysis/integration/notes/notes-improper}
